\section{Contexto anterior e recorte documental}
\claim{CL-C2-BOUNDARY}{A C1 já fornecia desempenho sustentado de 24 slots e ensaios curtos emparelhados de 24/32 slots. Não foi executado um novo controlo sustentado de 24 slots com o protocolo da C2. A comparação histórica entre campanhas não é, por isso, uma estimativa causal do efeito sustentado dos slots.}

\begin{tabularx}{\linewidth}{@{}p{37mm}Y@{}}
\toprule \textbf{Identidade} & \textbf{Pergunta a que responde}\\ \midrule
Experiência & Que modelo, binário, protocolo e workload produziram a observação?\\
Snapshot de evidência & Que revisão dos ficheiros publicados foi consultada?\\
Fontes do relatório & Que narrativa, seletores, classe e ferramentas geraram o documento?\\
Build & Que ambiente produziu estes bytes PDF?\\
\bottomrule
\end{tabularx}

\claim{CL-C2-IDENTITY}{O snapshot documental C2 é \Code{81531d9dc33c}. O resumo identifica \Code{c67529e29084} como base do laboratório e \Code{39f958c81eb8} como último commit de código medido. Um único SHA na capa não substituiria estas identidades.}

\claim{CL-C2-AUTHORITY}{O índice mantém o resumo inicial de C1 como histórico e identifica \Code{analysis-v2} para os detalhes corrigidos de eventos C2-C. Este relatório usa o gate e o resumo para as métricas centrais; não reexecuta a análise de eventos.}
